\chapter{Introduction}

Self-assembly pretains to spontaeneous formation of an organized structure from a pre-existing set of separate and disordered subunits/lower level components. The regularity, reversibility\cite{Reference13} and repeatability of these patterns indicate the existence of overarching laws/postulates that govern the overall interactions between the individual components. The postulates guide the motion of individual members which finally attain steady state by the balance of repulsive and attractive forces\cite{Reference11}.   

The existence of order and spontaeneous self-assembly has been of interest in the recent past\cite{Reference1}. In nature self-organisation has been observed to give rise to organized structures and spatial patterns both in physical and biological worlds\cite{Reference5}\cite{Reference6} e.g. planeraty weather patterns\cite{Reference11}, sand grains assembling to form dunes, lignification of cell walls in plants\cite{Reference2}\cite{Reference3}\cite{Reference4},in groups of microrganisms(bacteria), migration of fish\cite{Reference5}\cite{Reference6}, shape of honeycombs\cite{Reference7}\cite{Reference8}. Several attemps have been made to harness the structural order of materials and translate them into important practical applications. Examples being porous solids, materials with tunable mechanical properties\cite{Reference9}\cite{Reference10}, self-assembled monolayers(SAMs)\cite{Reference34}\cite{Reference35}.

Under the influence of flow-fields, many body suspensions with passive micro-particles have been known to show dynamic rearrangements leading to long-range ordered micro-structures e.g. microfluidic sorting devices, rouleaux formation in RBCs\cite{Reference12}\cite{Reference14}, advanced microstructured materials like microlens arrays\cite{Reference36} and nanostructured protein microfibres\cite{Reference37}.

In flow-fields under confinement microparticles have been known to interact via hydrodynamic interactions giving rise to long-range order. For instance, RBC suspensions in narrow capillaries\cite{Reference28}\cite{Reference29} form disjoined arrays of regular spacings. Several studies have reported the formation of droplet strings or trains and pearl necklace like morphologies in the suspensions of two immiscible polymers\cite{Reference15}\cite{Reference16}. Deformable particles especially under strong confinement readily rearrange in a wide range of volume fractions\cite{Reference30}\cite{Reference31}.

To develop a proper understanding of the phenomena, we had conducted studies aimed at investigating the confinement driven spontaeneous assembly of deformable particles under couette flow. Earlier we had suggested that the ordering of a  single layer of particles along the flow-direction resulted from the the hydrodynamic quadrupolar interactions between particles which assume a 2-D nature in the far-field\cite{Reference32}. Additionally, we had also suggested that while quadrupoles control the alignment of particles, the swapping-trajectory effect\cite{Reference27} controls their spacing albeit by the usage of a quasi-steady 2D model\cite{Reference33}. The present work specifically focusses on expanding the work to account for the three-dimensional dynamics of the system. 


%Self-assembly is the process of formation of an organized structure or pattern from pre-existing disordered subunits, the condition being the non-involvement of external factors, driven only by internal forces and interaction occurring within the system. Based on the nature of subunits involved, it can be termed as molecular, supramolecular, or nanoparticle self-assembly. The main driving forces can be the attainment of equilibrium or minimization of free energy or inter-unit interactions; the inter-unit interactions are mostly non-covalent in nature.

%The  central  tenet  of  self-organization  is  that  simple  repeated  interactions  between individuals can produce complex adaptive patterns at the level of the group.I argue that the key to understanding collective behaviour lies in identifying the principles of the behavioural algorithms followed by individual animals and of how information flows between the animals. These principles, such as positive feedback, response thresholds and individual integrity, are repeatedly observed in very different animal societies. There is a sense in which all these collective patterns are  regular  and  even  predictable. The submergence of the individual  brings  a  new  order  to  the  group.  Nevertheless,  understanding  how  coordinated  patterns emerge from a mass of interactions between individuals poses a difficult problem. The regularity of collective animal behaviour leaves us feeling that there must be some  unifying  laws  which  govern  these  different phenomena.

%What these diverse systems hold in common is the proximate means by which they acquire order and structure. In self-organizing systems, pattern at the global level emerges solely from interactions among lower-level components. Remarkably, even very complex structures result from the iteration of surprisingly simple behaviors performed by individuals relying on only local information.

%Here, we limit the term to process-es that involve pre-existing components (sep-arate  or  distinct  parts  of  a  disordered  struc-ture), are reversible, and can be controlled byproper  design  of  the  components.
%Their  steady-state positions balance attractions and repulsions. 
